%==============================================================================
\part{Verdicts}
%==============================================================================

\section{The verdict algebra}
\label{sec:verdicts}

\subsection{Three independent predicates}

Before naming the verdicts, we isolate what a verdict is about. Fix a step
$\Step = (x, \Src, \rho, \beta, b)$ executed against a source holding dataset
$D$, with the values of $\beta$ already computed. Three predicates apply, and
the burden of this subsection is that they are logically independent.

\begin{definition}[The three predicates]
\label{def:three}
\leavevmode
\begin{itemize}[leftmargin=2.4em]
\item $\mathrm{Exp}(\Step)$ holds iff the abstract request $\rho$ is
expressible against $\Src$: formally, $\Req(\rho) \subseteq \Capset(\Src)$.
\item $\mathrm{Der}(\Step)$ holds iff
$\llbracket \rho \rrbracket_D \neq \emptyset$ on the dataset the source in fact
holds, with the supplied bindings $\beta$.
\item $\mathrm{Ans}_b(\Step)$ holds iff the concrete request completes within
the step's budget $b$.
\end{itemize}
\end{definition}

\begin{proposition}[Independence]
\label{prop:independence}
The three predicates of \cref{def:three} are logically independent: all eight
truth assignments are realisable by some (step, source, dataset, budget).
\end{proposition}

\begin{proof}
We exhibit the constructions; each varies one predicate holding the others
fixed, and the eight cases follow by composing them.

To vary $\mathrm{Exp}$ alone: take a request using \textsf{regex} and swap the
source between one declaring \textsf{regex} and one not. Neither the dataset
nor the budget changes, so $\mathrm{Der}$ and $\mathrm{Ans}_b$ are untouched on
the branch where the request is issued at all; on the branch where it is not,
$\mathrm{Der}$ is still a property of $D$ and $\rho$, both fixed, and
$\mathrm{Ans}_b$ is vacuously assigned by convention.

To vary $\mathrm{Der}$ alone: fix the source and request and vary the dataset
between one containing a witness for $\phi$ and one not. Capability
declarations do not depend on $D$, so $\mathrm{Exp}$ is unchanged; a request
that completes on one finite dataset completes on the other for a budget
exceeding both costs, so $\mathrm{Ans}_b$ can be held.

To vary $\mathrm{Ans}_b$ alone: fix source, request and dataset, and vary $b$
across the cost of the request. Neither $\Capset(\Src)$ nor
$\llbracket\rho\rrbracket_D$ depends on $b$, so the other two are unchanged.

Since each predicate is separately variable with the other two held fixed, and
each is two-valued, all eight assignments are realisable.
\end{proof}

\Cref{prop:independence} is the reason a single boolean cannot carry the
outcome of a step: it must report a point in an eight-element space, and one
bit indexes two. The verdict set below is the coarsening of that space which
distinguishes the cases calling for different repairs.

\begin{remark}[Only one predicate carries a clock]
\label{rem:clock}
$\mathrm{Exp}$ is a property of the request and the declaration;
$\mathrm{Der}$ is a property of the request and the dataset; only
$\mathrm{Ans}_b$ mentions the budget. A system reporting a single outcome
cannot say which of the three it means, and a practitioner who reruns a failing
step with more time is implicitly betting that the failure was in the third.
The bet is unfounded two-thirds of the time, and \cref{lst:script} provides no
way to check it.
\end{remark}

\subsection{Six verdicts}

\begin{definition}[Verdict set]
\label{def:verdicts}
$V = \{\vans, \vempty, \vsurf, \vtime, \vref, \vstarve\}$, assigned to a step
by the following rules, applied in order:
\begin{enumerate}[label=(R\arabic*),leftmargin=3em]
\item If some $y \in \beta$ has a verdict other than $\vans$, or has verdict
$\vans$ with a payload whose retention fell below the step's declared
expectation, then $\vstarve$; the diagnosis names the offending predecessor.
\item Else if $\Req(\rho) \not\subseteq \Capset(\Src)$ then $\vsurf$; the
diagnosis names $\Req(\rho)\setminus\Capset(\Src)$.
\item Else if the plan's remaining budget is below $\cost_\Src$ evaluated at
the input cardinality, then $\vref$; the diagnosis names the shortfall.
\item Else issue the request. If it does not complete within $b$ then $\vtime$;
the diagnosis names $b$ and the elapsed cost.
\item Else if the extracted result set is empty then $\vempty$.
\item Else $\vans$, with the result set as payload.
\end{enumerate}
\end{definition}

The ordering of (R1)--(R6) is not arbitrary and is not merely an implementation
convenience: it is what makes the verdicts \emph{attributive}. Each rule fires
only when every earlier cause has been excluded, so a step reporting $\vempty$
has established that its inputs were adequate, its capabilities sufficient, its
budget available and its request completed --- and therefore that the emptiness
is a fact about the data rather than about the machinery. No rule below (R6)
can say that.

\begin{definition}[Blocker]
\label{def:blocker}
The \emph{blocker} of a non-$\vans$ verdict is
\[
\mathrm{blk}(\vsurf) = \bmdl,\quad
\mathrm{blk}(\vtime) = \beng,\quad
\mathrm{blk}(\vref) = \bbud,\quad
\mathrm{blk}(\vstarve) = \bcorp,
\]
and $\vempty$ has no blocker.
\end{definition}

\begin{remark}[$\vempty$ has no blocker on purpose]
Assigning a blocker to $\vempty$ would assert that something obstructed the
step, and (R6) fires exactly when nothing did. A verdict of $\vempty$ is a
successful measurement of an empty set. Conflating it with an obstruction is
defect (D2) of \cref{sec:intro} in its purest form, and the discipline that
prevents it is that the blocker map is partial.
\end{remark}

\subsection{Distinguishability, and the one-bit interface}

\begin{theorem}[Six-way distinguishability]
\label{thm:six}
For each pair $v \neq v'$ in $V$ there exist a plan, two source
configurations agreeing on everything the plan can observe except as noted, and
a step whose verdict is $v$ in the first and $v'$ in the second, such that the
executor of \cref{def:verdicts} reports the difference. Consequently the six
verdicts are pairwise distinguishable by a plan executor.
\end{theorem}

\begin{proof}
It suffices to realise each verdict by a configuration differing from a common
baseline in exactly one respect, since then any two verdicts are separated by
the pair of configurations realising them.

Take as baseline a step whose predecessors all returned $\vans$ at full
retention, whose source declares every capability the request uses, whose
budget exceeds the request cost, and whose dataset contains a witness. By
(R1)--(R6) the baseline verdict is $\vans$.

$\vempty$: remove the witness from $D$, changing nothing else. (R1)--(R4) still
pass; (R5) fires.

$\vtime$: restore the witness and set $b$ below the request cost. (R1)--(R3)
pass, (R4) fires.

$\vref$: restore $b$ and set the plan's remaining budget below
$\cost_\Src$. (R1)--(R2) pass, (R3) fires.

$\vsurf$: restore the budget and remove one capability of $\Req(\rho)$ from
$\Capset(\Src)$. (R1) passes, (R2) fires.

$\vstarve$: restore the capability and set a predecessor's verdict to any
non-$\vans$ value, or its retention below the declared expectation. (R1) fires.

Each configuration differs from the baseline in exactly one respect and yields
a distinct verdict, and the executor computes the verdict by the stated rules,
so it reports each. The six are therefore pairwise distinguishable.
\end{proof}

\begin{corollary}[A row-returning interface conflates at least four verdicts]
\label{cor:onebit}
Let $I$ be an interface whose response to a request is a result set, with
failure signalled by returning the empty set. Then $I$ cannot distinguish
$\vempty$, $\vtime$, $\vsurf$ and $\vstarve$ in any of the four configurations
constructed in the proof of \cref{thm:six}. No discipline imposed on callers of
$I$ repairs this.
\end{corollary}

\begin{proof}
In each of the four configurations the value $I$ returns is the empty result
set: in the $\vempty$ configuration because the denotation is empty; in the
$\vtime$ configuration because the request did not complete and $I$ has no
other channel; in the $\vsurf$ configuration because the request was not issued
or was issued and degraded; in the $\vstarve$ configuration because the inputs
were empty and a request over empty inputs has empty denotation. The four
responses are equal as values. A caller's behaviour is a function of the values
it receives, so no caller distinguishes them, and no discipline on callers
changes the values. Hence the conflation is a property of $I$, not of its use.
\end{proof}

\Cref{cor:onebit} is the formal content of defect (D2), and it is worth being
explicit about its scope. It does not say that existing systems are badly
engineered. It says that the interface contract they expose --- rows, with
emptiness overloaded as failure --- has insufficient capacity, and that a
system built on it cannot recover the distinction downstream however carefully
it is written. The repair is not better error handling; it is a wider return
type.

\subsection{Starvation is the characteristic federated failure}

The verdict $\vstarve$ deserves separate treatment, because it is the one with
no single-database analogue, and because it is not a corner case.

\begin{proposition}[Starvation is unreachable in the closed setting and reachable in the federated one]
\label{prop:starve-reachable}
Consider the restriction of \cref{def:verdicts} to plans of a single step
against a fixed corpus --- the closed setting. Then:
\begin{enumerate}[label=(\alph*),leftmargin=2.4em]
\item no execution assigns $\vstarve$, and consequently the blocker $\bcorp$
is emitted by no rule;
\item in the federated setting with $m \geq 2$ steps, $\vstarve$ is assigned on
a set of configurations of positive measure in the following sense: for every
step $\Step_i$ with $\beta_i \neq ()$ and every predecessor $\Step_j$ feeding
it, there is a dataset perturbation at $\Src_j$ alone that causes $\Step_i$ to
report $\vstarve$ while $\Step_i$'s own source, capabilities and budget are
untouched.
\end{enumerate}
\end{proposition}

\begin{proof}
(a) In a single-step plan, $\beta = ()$, so the antecedent of (R1) is vacuously
false and (R1) never fires. Since (R1) is the unique rule assigning $\vstarve$,
the verdict is unreachable; and since $\bcorp$ is by \cref{def:blocker} the
image of $\vstarve$ alone, no blocker $\bcorp$ is emitted. The corpus is fixed
and given, so its inadequacy manifests as $\vempty$ under (R5) --- an empty
answer, correctly reported as such, with no obstruction named.

(b) Let $\Step_j \to \Step_i$ be an edge of $G(\Plan)$ with $y \in \beta_i$
bound by $\Step_j$. Perturb the dataset at $\Src_j$ by deleting the witnesses
for $\rho_j$. Then $\Step_j$ reports $\vempty$ by (R5); its capability set,
budget and lowering are unchanged, so no other verdict applies to it. At
$\Step_i$, the antecedent of (R1) now holds, since $y$ has a verdict other than
$\vans$; (R1) fires and $\Step_i$ reports $\vstarve$ naming $\Step_j$. Nothing
at $\Src_i$ was altered.
\end{proof}

\begin{remark}[Why this is the characteristic case]
\label{rem:starve-characteristic}
In the closed setting the corpus is an input: if it lacks the assertions a
question needs, that is a fact about the world one reports and does not
diagnose. In federation the corpus of step $i$ is not an input but an
\emph{artefact of steps $1..i-1$}, and its inadequacy is therefore a
diagnosable event with a locatable cause. \Cref{prop:starve-reachable} says
that federation makes a previously unreachable diagnosis reachable, and
\cref{prop:starve-common} below says it is the modal outcome once translation
enters. Read together with \cref{cor:onebit}: the failure mode that federation
newly creates is exactly one of the four that the conventional interface cannot
see.
\end{remark}

\begin{proposition}[Starvation dominates once translation enters]
\label{prop:starve-common}
Let $\Plan$ contain a translation step whose retention is $r < 1$
(\cref{def:retention}), feeding a step whose per-input hit probability is $p$
independently across inputs. Then for input cardinality $N$ the probability
that the downstream step returns a non-empty payload is
$1 - (1-p)^{rN}$, whereas absent the translation it is $1-(1-p)^{N}$. For $p$
small and $rN$ small the ratio of failure probabilities approaches
$e^{-prN}/e^{-pN} = e^{p N (1-r)}$, so the excess failure attributable to
translation grows exponentially in the lost fraction $1-r$ times the input
size.
\end{proposition}

\begin{proof}
The translation delivers $rN$ inputs by \cref{def:retention}. Under
independence the downstream step misses on all of them with probability
$(1-p)^{rN}$, giving the stated success probability; the untranslated case is
$r=1$. Taking the ratio of the two failure probabilities and using
$(1-p)^{k} = e^{k\ln(1-p)} \approx e^{-pk}$ for small $p$ yields
$e^{-prN + pN} = e^{pN(1-r)}$.
\end{proof}

The independence hypothesis in \cref{prop:starve-common} is false in practice
--- identifiers that fail to translate are correlated with identifiers that are
poorly annotated downstream, so the true excess is larger than the bound and
not smaller. We state the independent case because it is the one we can prove,
and record the direction of the bias rather than claiming the bound is tight.

\subsection{Verdict propagation}

\begin{definition}[Propagation]
\label{def:propagation}
A plan's verdict is $\vans$ if every emitted step is $\vans$; otherwise it is
the verdict of the earliest non-$\vans$ step in the sequence order, together
with the transitive closure of blame along $G(\Plan)$.
\end{definition}

\begin{proposition}[Blame is well defined and terminates at a non-starved step]
\label{prop:blame}
For any plan and execution, following the diagnosis of a $\vstarve$ verdict
from step to named predecessor terminates, and terminates at a step whose
verdict is not $\vstarve$.
\end{proposition}

\begin{proof}
By \cref{def:plan}, every variable in $\beta_i$ is bound by some $\Step_j$ with
$j < i$, so the named predecessor is strictly earlier in the sequence order.
The chain of blame is therefore strictly decreasing in a well-founded order on
a finite set and must terminate. It terminates at a step which is not
$\vstarve$, because a $\vstarve$ step always names a predecessor by (R1) and so
is never terminal; and the first step of the plan has $\beta_1 = ()$, so (R1)
cannot fire on it.
\end{proof}

\Cref{prop:blame} is what makes a $\vstarve$ report actionable rather than
circular: the analyst follows the chain and arrives at a step whose failure is
attributed to a capability, a budget, a timeout, or a genuinely empty dataset
--- one of the four things one can actually do something about. This is defect
(D3) resolved, and it is resolved by the structure of the verdict algebra
rather than by a convention that callers are asked to observe.
